\chapter{L'apparato circolatorio}

% ---------------------------------------------------------------------------
\section{Generalità}

\textbf{Cuore}: pompa. \textbf{Vasi sanguiferi}: vie di conduzione del sangue.

L'apparato circolatorio è l'insieme di organi deputati al trasporto dei liquidi: è un sistema
chiuso di vasi, in cui il sangue circola sotto la spinta del cuore.

\rubrica{Caratteristiche dei vasi sanguigni}
\begin{itemize}
  \item \textbf{resistenti}, per sopportare le variazioni della pressione ematica, che varia con
    la distanza dal cuore;
  \item \textbf{flessibili}, per adattarsi ai movimenti degli organi vicini;
  \item \textbf{robusti}, per evitare di essere occlusi da pressioni esterne;
  \item \textbf{adattati nella struttura} alla specifica funzione.
\end{itemize}

I vasi sanguigni si suddividono in: arterie, arteriole, capillari, venule, vene.

\rubrica{Funzioni dei vasi}
\begin{itemize}
  \item trasporto del sangue;
  \item scambio gassoso anidride carbonica/ossigeno nei capillari polmonari, e scambio gassoso
    ossigeno/anidride carbonica e di nutrienti/sostanze di rifiuto nei capillari sistemici;
  \item regolazione della pressione sanguigna;
  \item indirizzamento del flusso sanguigno sistemico verso i distretti che in quel dato istante
    necessitano di un maggiore apporto di sangue: es.\ l'apparato digerente durante la digestione,
    o i muscoli durante l'attività fisica.
\end{itemize}

% ---------------------------------------------------------------------------
\section{Struttura istologica dei vasi sanguigni}

\begin{itemize}
  \item \textbf{tonaca intima}: endotelio monostratificato piatto e sottostante strato di tessuto
    elastico (endotelio + strato subendoteliale fibroso, o membrana basale, + membrana elastica
    interna; quest'ultima più spessa nelle arterie);
  \item \textbf{tonaca media}: strati concentrici di tessuto muscolare liscio con fibre collagene
    (nelle arterie è presente anche una banda esterna di tessuto elastico, la membrana elastica
    esterna);
  \item \textbf{tonaca esterna} (o avventizia): tubo fibroso robusto contenente vasi (\textit{vasa
    vasorum}) e nervi, che ancora i vasi al connettivo circostante.
\end{itemize}

% ---------------------------------------------------------------------------
\section{L'albero arterioso}

L'albero arterioso si struttura in:
\begin{enumerate}
  \item \textbf{arterie}, in cui è fondamentale la componente elastica;
  \item \textbf{arteriole}, in cui è fondamentale la componente muscolare;
  \item \textbf{capillari arteriosi};
\end{enumerate}
a questi fanno seguito: capillari venosi, venule post-capillari, vene.

\rubrica{Arterie elastiche}
Sono le più vicine al cuore, in cui prevale la componente elastica; ricevono una forte spinta
sistolica, ed hanno un calibro elevato ($1$--$3\,cm$). Esempio: l'\textbf{aorta}. Contengono una
notevole quantità di elastina: durante la sistole ventricolare si espandono incamerando energia
potenziale, che restituiscono durante la diastole ventricolare, contribuendo a:
\begin{itemize}
  \item spingere il sangue in avanti;
  \item smorzare le oscillazioni pressorie;
  \item trasformare il flusso da intermittente a continuo.
\end{itemize}

\rubrica{Arterie muscolari}
Prevale la componente muscolare; diametro $0,5\,mm$--$1\,cm$. Esempio: l'\textit{arteria
femorale} (spessa lamina elastica).

\rubrica{Arteriole}
Scompare l'avventizia, la tonaca media è muscolare; diametro $< 0,5\,mm$.

% ---------------------------------------------------------------------------
\section{Capillari}

\textbf{Capillari}: cellule endoteliali + lamina basale; diametro $5$--$10\,\mu m$; lunghezza
media $1\,mm$. Sono i vasi più piccoli, deputati allo scambio coi tessuti: spesso una sola cellula
endoteliale forma tutta la sezione trasversa del capillare.

Sono molto numerosi in tutti i tessuti, con alcune eccezioni: tendini e legamenti; cartilagine;
epiteli; cornea.

Essendo molto numerosi ed avendo una parete molto sottile, rappresentano le strutture ideali per
fornire alle cellule sostanze indispensabili quali gas, nutrienti, ormoni.

\rubrica{Tipi di capillari}
\begin{itemize}
  \item \textbf{continui}: i più comuni nell'organismo, abbondanti nella pelle e nei muscoli. Le
    cellule endoteliali formano un rivestimento continuo; giunzioni serrate mantengono le cellule
    aderenti, lasciando solo piccoli spazi dai quali escono piccole quantità di liquidi;
  \item \textbf{fenestrati}: simili ai continui, eccetto per la presenza di pori che attraversano
    a tutto spessore la cellula. Questi pori sono coperti da sottili diaframmi, e poi dalla lamina
    basale delle cellule endoteliali; le fenestrature consentono un notevole aumento di
    permeabilità. Si trovano preferenzialmente dove devono avvenire in maniera rilevante:
    \begin{itemize}
      \item formazione di ultra-filtrato (rene, plessi corioidei);
      \item riassorbimento di liquidi e sostanze in esso disciolte (digerente, rene);
      \item contatti con l'apparato vascolare: organi endocrini;
    \end{itemize}
  \item \textbf{sinusoidi}: simili ai fenestrati, ma con forma irregolare e tortuosa; pori
    ingranditi, spazi anche ampi tra le cellule endoteliali; lamina basale più sottile; lento
    movimento del sangue. Queste modificazioni consentono sia ai liquidi che a macromolecole e
    cellule del sangue di entrare nello spazio interstiziale o nel lume del capillare. Presenti in
    fegato, milza, midollo osseo ed organi endocrini: con interruzioni intercellulari endoteliali
    e membrana basale incompleta, permettono anche il passaggio di cellule.
\end{itemize}

\rubrica{Edema}
Ristagno di liquido nei tessuti. Cause:
\begin{itemize}
  \item aumento della pressione venosa (insufficienza cardiaca o ristagno venoso);
  \item ristagno del flusso linfatico;
  \item capillari divenuti permeabili alle proteine plasmatiche (infiammazione, risposte
    immunitarie);
  \item diminuzione delle proteine plasmatiche (malnutrizione, malattie epatiche, etc.: ascite).
\end{itemize}

% ---------------------------------------------------------------------------
\section{Vene}

Struttura analoga alle arterie ma con minore componente muscolare ed elastica, e più connettivo.
Anche le vene sono innervate dall'ortosimpatico, che facendole contrarre costringe il sangue a
fuoriuscire da alcuni organi (fegato, milza, polmoni, ad elevata funzione capacitativa) per
incrementare la pressione del sangue.

Dai capillari il sangue passa nelle: venule post-capillari ($10$--$25\,\mu m$)
$\rightarrow$ venule collettrici ($20$--$50\,\mu m$) $\rightarrow$ vene muscolari
($50$--$100\,\mu m$) $\rightarrow$ vene (fino a $4\,cm$).

\rubrica{Confronto vene/arterie}
\begin{itemize}
  \item \textbf{vene}: parete più sottile; lume collassabile; struttura scarsamente elastica;
    possono essere dotate di valvole; bassa pressione;
  \item \textbf{arterie}: parete più spessa; forma cilindrica sempre conservata; struttura molto
    elastica; prive di valvole; alta pressione.
\end{itemize}

Nella maggior parte dei casi le vene hanno un decorso speculare a quello delle arterie, di cui
assumono il nome e che accompagnano nel decorso.

Le vene degli arti e della parte inferiore del corpo presentano \textbf{valvole}, pieghe della
tonaca intima coperte da endotelio, che ricordano e funzionano in maniera simile alle valvole
semilunari.

\rubrica{Fistola artero-venosa}
Comunicazione anomala tra un'arteria e una vena: lesione a carico di un'arteria e di una vena
affiancate. La lesione è solitamente rappresentata da una ferita penetrante, come quella provocata
da un coltello o un proiettile. Può essere congenita (rara) o acquisita.

% ---------------------------------------------------------------------------
\section{Capillari e plessi capillari}

I capillari costituiscono l'unico distretto cardiovascolare ove sono possibili scambi tra sangue
e liquido interstiziale circostante, e quindi cellule e tessuti. Costituiscono una rete a maglie
intricate (plesso capillare, o letto capillare) in cui il sangue scorre a velocità estremamente
bassa, utile per la diffusione o il trasporto attivo transmembrana dei materiali.

Normalmente il sangue procede dai capillari alle vene, per poi far ritorno al cuore. Vi sono 3
eccezioni:
\begin{itemize}
  \item \textbf{ipofisi}: circolazione portale ipotalamo-ipofisaria;
  \item \textbf{fegato}: circolazione portale;
  \item \textbf{rene}: il nefrone presenta un letto capillare arterioso seguito da uno
    arterovenoso.
\end{itemize}

% ---------------------------------------------------------------------------
\section{Circolazione}

\rubrica{La grande circolazione}
Detta anche \textbf{circolazione sistemica}: è la circolazione sanguigna tra il cuore, i tessuti
dell'organismo e di ritorno al cuore. Dal cuore partono le arterie, nelle quali circola il sangue
ricco di ossigeno che raggiunge i tessuti dell'organismo, dove verrà ceduto l'ossigeno ed in
cambio verrà acquistata anidride carbonica. Dai tessuti, il sangue ricco di anidride carbonica
farà ritorno al cuore attraverso le vene.

\rubrica{La piccola circolazione}
Detta anche \textbf{circolazione polmonare}: è la circolazione sanguigna tra il cuore, i polmoni e
di ritorno al cuore. Il sangue ricco di anidride carbonica proveniente dalla circolazione
sistemica viene pompato dal cuore verso i polmoni, dove avviene lo scambio gassoso: viene ceduta
anidride carbonica ed acquistato ossigeno. Il sangue ritorna al cuore, il quale lo pompa nella
circolazione sistemica verso i tessuti dell'organismo.

Il cuore è quindi una pompa formata funzionalmente da 2 metà:
\begin{itemize}
  \item \textbf{destra}, che riceve e spinge sangue deossigenato;
  \item \textbf{sinistra}, che riceve e spinge sangue ossigenato.
\end{itemize}

I ventricoli conferiscono al sangue una consistente pressione: il sistema circolatorio è quindi
organizzato in 2 componenti, una ad alta pressione in partenza dal cuore, e una a bassa pressione
in arrivo al cuore.

% ---------------------------------------------------------------------------
\section{Aorta}

L'aorta si distingue in 3 porzioni: aorta ascendente (dalla quale originano le coronarie), arco
dell'aorta (dal quale originano il tronco brachiocefalico, la carotide comune e la succlavia di
sinistra), aorta discendente (a sua volta suddivisa in 2 porzioni consecutive: toracica ed
addominale). Termina a livello di L4, dividendosi nelle 2 arterie iliache comuni, destra e
sinistra.

\rubrica{Aorta ascendente}
Inizia dalla valvola semilunare aortica, all'altezza della 3\textsuperscript{a} cartilagine
costale di sinistra, margine inferiore. All'origine sono presenti i 3 seni aortici (di Valsalva),
quindi un rigonfiamento detto \textbf{bulbo aortico}. Si porta in alto e a destra, e all'altezza
del margine superiore della 2\textsuperscript{a} cartilagine costale trapassa nell'arco aortico.
Rami: le 2 arterie coronarie, destra e sinistra.

\rubrica{Arco aortico}
Si porta posteriormente, scavalca la biforcazione del tronco polmonare, passa a sinistra della
trachea e poi dell'esofago, passa superiormente al bronco di sinistra, e raggiunge il corpo di
T4. Genera l'arteria anonima (o tronco brachiocefalico), la carotide comune di sinistra e la
succlavia di sinistra. Nel feto, l'arteria polmonare comunica con l'arco aortico mediante il
\textit{condotto di Botallo}.

\rubrica{Aorta discendente}
Porta il sangue direttamente dal cuore al resto del corpo. Si distinguono 2 porzioni: toracica e
addominale; il limite è dato dal passaggio nel diaframma, a livello della dodicesima vertebra
toracica. Da entrambe le porzioni originano rami collaterali che irrorano le pareti (rami
parietali) o i visceri (rami viscerali).

\rubrica{Aorta discendente toracica}
\begin{itemize}
  \item arterie \textbf{bronchiali}: servono bronchi e polmoni;
  \item arterie \textbf{pericardiche}: servono il pericardio;
  \item arterie \textbf{mediastiniche}: servono gli organi mediastinici;
  \item arterie \textbf{esofagee}: servono l'esofago;
  \item arterie \textbf{freniche superiori}: servono il diaframma;
  \item arterie \textbf{intercostali}: servono la parete toracica.
\end{itemize}

\rubrica{Aorta addominale}
Porzione inferiore dell'aorta discendente. Genera:
\begin{itemize}
  \item rami \textbf{posteriori}, parietali, pari;
  \item rami \textbf{laterali}, viscerali, pari (reni, surreni, genitali);
  \item rami \textbf{anteriori}, viscerali, impari.
\end{itemize}

Nel dettaglio:
\begin{itemize}
  \item rami viscerali \textbf{anteriori}:
  \begin{enumerate}
    \item \textbf{tronco celiaco}: 3 rami, al fegato, alla cistifellea, all'esofago, allo stomaco,
      al duodeno, al pancreas, alla milza;
    \item \textbf{mesenterica superiore}: al pancreas e al duodeno, al piccolo intestino e al
      colon;
    \item \textbf{mesenterica inferiore}: al colon terminale e al retto;
  \end{enumerate}
  \item rami viscerali \textbf{laterali}: surrenali (alle ghiandole surrenali), renali (ai reni),
    spermatiche (alle gonadi: ovaie e testicoli);
  \item rami \textbf{parietali}: lombari (al dorso).
\end{itemize}
